\documentclass[11pt]{article}
\usepackage[margin=1in]{geometry}
\usepackage{amsmath,amssymb,amsthm}
\usepackage{booktabs}
\usepackage{graphicx}
\usepackage{hyperref}
\usepackage{xcolor}

\newtheorem{theorem}{Theorem}
\newtheorem{proposition}{Proposition}
\newcommand{\stack}{\textsc{Stack}}
\newcommand{\auto}{\textsc{AutoSelect}}
\newcommand{\relg}{\textsc{RelGate}}

\title{\textbf{When Does Reliability Help Universal Anomaly Detection?\\
A Strong Stacking Baseline, a Rigorous Negative Result, and a Mechanism}}
\author{}
\date{}

\begin{document}
\maketitle

\begin{abstract}
No single anomaly detector dominates across domains, which has motivated a wave of
\emph{reliability-aware} meta-routers that select, fuse, or abstain among detector
experts using validation quality, score-distribution drift, inter-detector
disagreement, and calibration signals. We ask a blunt question: does any of this
machinery beat the obvious baseline of picking, or stacking, detectors by
validation performance? On a broad benchmark of \textbf{123 anomaly-detection tasks
across five families} (tabular, image-OOD, text, cyber, fraud) we report three
results. \textbf{(1) Positive:} a parameter-light rank-normalized logistic
\emph{stack} is a strong, underused universal baseline---it beats per-task
validation-AUROC selection by $+0.036$ AUROC (95\% CI $[0.023,0.050]$) and the best
fixed detector by $+0.075$ ($[0.052,0.101]$), and even exceeds the best-single-detector
oracle. \textbf{(2) Negative:} reliability-aware routing does \emph{not} beat these
simple baselines. Across three deployment-degradation regimes (i.i.d., uniform shift,
heterogeneous missingness) and two independent router designs, with clean ablations,
reliability features carry no significant signal beyond validation quality and
detector identity; an explicit reliability gate on the stack \emph{hurts} ($-0.037$,
$[-0.050,-0.026]$). \textbf{(3) Mechanism:} we show in a controlled simulation that
reliability gating helps \emph{if and only if} a channel can fail independently while
others stay clean. Single-modality detectors share one input, so degradation is
shared and the condition is violated---explaining the negative result---whereas
genuine multimodal sensing satisfies it. We pre-register a falsifiable multimodal
test for exactly that regime. Our contribution is an honest map of where
reliability-aware anomaly routing helps, where it does not, and why.
\end{abstract}

\section{Introduction}
Anomaly detection spans tabular records, industrial vision, RGB-D, time series,
cyber traffic, fraud, and out-of-distribution image classification. A robust
empirical regularity, confirmed again here, is that \emph{no single detector wins
everywhere}: the best detector on one task is mediocre on the next (Table~\ref{tab:main},
fixed detectors span ranks $5.4$--$8.7$ of $11$). This non-dominance is the practical
driver of \emph{negative transfer}---deploying a detector chosen on one domain and
losing on another.

The natural response is to choose per task. The simplest choice rule, picking the
detector with best validation AUROC (\auto), is already strong and is the honest
baseline any meta-method must beat. A large recent literature instead proposes
\emph{reliability-aware} routers that estimate, at deployment, which experts to
trust---via score-distribution drift, inter-detector disagreement, calibration, and
sharpness---and select, fuse, or abstain accordingly. The intuition is compelling:
when validation performance becomes \emph{stale} under deployment shift, a reliability
signal computed at test time should rescue the routing decision.

We test that intuition as hard as we can and find it does not hold on real
single-modality benchmarks. Our contributions:
\begin{enumerate}
\item \textbf{A strong baseline.} A rank-normalized logistic stack beats validation
selection and every fixed detector across 123 tasks / 5 families, with no test-label
leakage (\S\ref{sec:results}). It should be the default comparator for universal
anomaly detection; it is rarely reported.
\item \textbf{A rigorous negative result.} Reliability-aware routing does not beat it.
We show this for two router families (a learned gradient-boosted selector and a
hand-designed gate/stack) across three degradation regimes, each with a
reliability-features-removed ablation that comes back null or negative
(\S\ref{sec:negative}). We also document a parallel pipeline whose ``ablation''
silently held all variants identical---a cautionary failure mode.
\item \textbf{A mechanism.} A controlled simulation isolates the single structural
variable that decides the matter: reliability gating helps iff channels fail
\emph{independently}. Single-modality detectors share one input, violating it
(\S\ref{sec:mechanism}). This reconciles the negative result with the operating-boundary
theory (\S\ref{sec:theory}) and motivates a pre-registered multimodal test
(\S\ref{sec:multimodal}).
\end{enumerate}
We deliberately do not claim a universal state-of-the-art detector. We claim an honest
characterization of when a popular idea works.

\section{Related Work}
\textbf{Universal benchmarks and non-dominance.} Large suites such as ADBench
established that no detector dominates and that model selection matters; our 123-task
set extends this view across five families using pre-extracted feature
representations (ResNet for image-OOD, BERT for text). \textbf{Model selection /
meta-learning} (MetaOD and unsupervised outlier model selection) pick a detector per
dataset from meta-features; \auto{} is the strong label-aware analogue and our primary
baseline. \textbf{Fusion, stacking, and calibration} combine experts via rank mean,
confidence-weighted mean, or learned stacking; we find rank-normalized stacking is the
strongest of these and beats selection. \textbf{Drift detection and selective
prediction} provide the reliability signals (KS drift, disagreement, abstention) that
reliability-aware routers use; our negative result is precisely about whether those
signals add value to selection/fusion for anomaly detection.

\section{Problem Setup}
For task $t$ we have detectors $D_1,\dots,D_M$ producing scores $s_m(x)$. We split each
task into train (fit detectors, unlabeled), validation (labeled), and test (labeled,
held out). $A_V(m)$ and $A_T(m)$ denote validation and test AUROC; $m^\star=\arg\max_m
A_T(m)$ is the best-single-detector oracle (an upper bound for \emph{selection}, not
for fusion). \emph{Regret} is $A_T(m^\star)-A_T(h)$ for routed scores $h$.

\textbf{Leakage discipline.} No test labels are used to fit any router, stacker, gate,
or threshold. Validation labels are permitted (they are available at deploy time).
Unlabeled test-score \emph{distributions} may inform reliability features---this is
deployment monitoring, not label use---and we mark it as such.

\section{Methods}
We compare: each \textbf{fixed} detector; \auto{} (per-task $\arg\max_m A_V(m)$);
\textbf{rank\_mean} and \textbf{cw\_mean} (naive ensembles); \stack, a logistic
regression trained on \emph{rank-normalized validation scores} $r_m(x)=\mathrm{rank}_m(x)/n$
and applied to rank-normalized test scores; and reliability variants. The
\textbf{reliability gate} $\relg$ reweights channels by validation quality times an
unlabeled test-time consensus-agreement term; its \textbf{ablation} uses validation
quality only. A learned router (gradient-boosted regressor over detector identity,
validation AUROC, dataset meta-features, and---in the reliability arm---score
dispersion, gini, skew/kurtosis, inter-detector disagreement, and val$\to$test KS
drift) predicts per-detector test AUROC under leave-task-out and leave-family-out
cross-validation.

\section{Benchmark and Protocol}
123 tasks: tabular (43, ADBench Classical), image-OOD (61, ADBench CV/ResNet18), text
(13, ADBench NLP/BERT), cyber (5, UNSW-NB15 splits), fraud (1, credit-card). Detector
zoo: ECOD, COPOD, IForest, LOF, KNN, OCSVM. Each task: $50/25/25$ stratified
train/val/test, standardized on train. Metrics: per-task average rank, mean AUROC,
regret-to-oracle, expected calibration error (ECE), and paired bootstrap CIs over
tasks ($10^4$ resamples). Per-family results are reported to expose worst-family
behavior. All scores are persisted to a score archive so routing, ablation, and
calibration iterate without re-running detectors.

\section{Results: a strong stacking baseline}\label{sec:results}
Table~\ref{tab:main} reports the main comparison. Stacking is the best non-oracle
strategy and \emph{exceeds the best-single-detector oracle} ($0.784$ vs $0.756$):
fusion beats selection because diverse detectors are complementary. Validation
selection is itself strong---well above any fixed detector---confirming that the right
baseline for ``universal'' anomaly detection is per-task selection, not a fixed default.

\begin{table}[t]\centering
\caption{Main benchmark, 123 tasks / 5 families. Average rank is among the 11 non-oracle
strategies (lower is better). The oracle is the best \emph{single} detector per task
(a selection ceiling; fusion may exceed it). Leakage-free; see \S\ref{sec:results}.}
\label{tab:main}
\begin{tabular}{lrrrr}
\toprule
Strategy & Avg rank & Mean AUROC & Regret & ECE \\
\midrule
\stack{} (rank-norm logistic) & \textbf{2.53} & \textbf{0.784} & --- & 0.271 \\
\auto{} (val-AUROC) & 3.38 & 0.748 & 0.008 & 0.393 \\
\stack{}+reliability gate & 4.00 & 0.746 & 0.010 & 0.325 \\
fixed/KNN (best fixed) & 5.42 & 0.709 & 0.048 & --- \\
rank\_mean & 6.53 & 0.707 & 0.049 & 0.407 \\
cw\_mean & 6.66 & 0.708 & 0.048 & 0.396 \\
fixed/\{LOF,OCSVM,IForest,ECOD,COPOD\} & 6.4--8.7 & 0.68--0.70 & 0.05--0.08 & --- \\
\midrule
oracle (best single detector) & --- & 0.756 & 0.000 & --- \\
\bottomrule
\end{tabular}
\end{table}

\begin{table}[t]\centering
\caption{Key contrasts (paired bootstrap over 123 tasks, $10^4$ resamples).}
\label{tab:contrasts}
\begin{tabular}{lrrc}
\toprule
Contrast & $\Delta$AUROC & 95\% CI & Significant \\
\midrule
\stack{} $-$ \auto{}            & $+0.036$ & $[0.023,\,0.050]$ & yes \\
\stack{} $-$ best fixed         & $+0.075$ & $[0.052,\,0.101]$ & yes \\
\auto{} $-$ best fixed          & $+0.040$ & $[0.023,\,0.061]$ & yes \\
\stack{}+gate $-$ \stack{} (\emph{reliability ablation}) & $-0.037$ & $[-0.050,\,-0.026]$ & \textbf{hurts} \\
\bottomrule
\end{tabular}
\end{table}

Per-family average rank confirms the stack is robust, not driven by one family:
tabular $2.41$, image-OOD $2.49$, cyber $1.00$, text $3.69$; it loses only on fraud
($n{=}1$, underpowered). Crucially, adding a reliability gate to the stack
\emph{reduces} AUROC by $0.037$ (Table~\ref{tab:contrasts})---our first and cleanest
sign that reliability signals do not help here.

\section{Results: reliability-aware routing does not help}\label{sec:negative}
We test the reliability premise directly, twice, with clean ablations
(Table~\ref{tab:negative}). A learned selector with the full reliability-feature
battery is \emph{worse} than plain validation selection on i.i.d.\ data ($-0.005$, CI
excludes $0$), and the reliability-vs-no-reliability ablation is null ($+0.001$, CI
$[-0.005,0.006]$). Under uniform deployment shift a learned router does beat \emph{stale}
selection, but a leave-one-out attribution shows the gain comes from \emph{detector-identity
priors} (which architectures survive corruption), not reliability features: the
reliability ablation is null at every severity. Under heterogeneous per-task
missingness---the regime most favorable to reliability---the ablation flickers around
zero non-monotonically ($-0.003,+0.006,-0.003,+0.015$) and never reaches significance.

\begin{table}[t]\centering
\caption{Reliability-feature ablations. $\Delta = $ (with reliability) $-$ (without);
$>0$ means reliability helps. None is significant; the stacker gate is significantly
negative. Two router designs, three regimes.}
\label{tab:negative}
\begin{tabular}{llrc}
\toprule
Regime / design & comparison & $\Delta$AUROC & CI excludes 0? \\
\midrule
i.i.d., learned selector & rel.\ vs no-rel.\ & $+0.001$ & no \\
i.i.d., learned selector & rel.\ vs \auto{} & $-0.005$ & yes (worse) \\
uniform shift (sev.\ 3), learned & rel.\ vs no-rel.\ & $+0.002$ & no \\
heterog.\ missingness (frac.\ 0.7), learned & rel.\ vs no-rel.\ & $+0.015$ & no \\
123-task benchmark, stacker gate & gate vs \stack{} & $-0.037$ & yes (worse) \\
\bottomrule
\end{tabular}
\end{table}

\textbf{A cautionary note on ablations.} A parallel pipeline reported a polished
``reliability ablation'' in which \emph{every} variant (full, no-reliability, no-drift,
no-disagreement, no-calibration) returned byte-identical metrics, because all were
silently mapped to the same model---the reliability features were never in the method.
We flag this because an ablation that does not ablate is indistinguishable, at a glance,
from one that shows components ``all contribute,'' and only inspection of identical
digits reveals it.

\section{Mechanism: reliability needs an independent clean channel}\label{sec:mechanism}
Why does a structurally reasonable idea fail? We isolate the cause in a controlled
simulation (Table~\ref{tab:mechanism}) with two regimes differing only in how
degradation acts. In \textbf{SHARED}, all channels are weakened together (the
single-modality analogue: one input feeds all detectors); in \textbf{INDEPENDENT}, one
channel loses its signal while the others stay clean (the genuine multimodal case). The
reliability gate, validation selection, and ablation are identical across regimes.

\begin{table}[t]\centering
\caption{Mechanism simulation (300 trials/regime, paired bootstrap). Reliability gating
helps iff a channel fails independently.}
\label{tab:mechanism}
\begin{tabular}{lrrr}
\toprule
 & \relg{} $-$ mean & \relg{} $-$ \auto{} & \relg{} $-$ ablation \\
\midrule
SHARED (all weaken together) & $-0.016$ (ns) & $+0.081$ & $-0.019$ (ns) \\
INDEPENDENT (one fails, rest clean) & $+0.026$ \checkmark & $+0.144$ \checkmark & $+0.024$ \checkmark \\
\bottomrule
\end{tabular}
\end{table}

Under SHARED degradation reliability gating cannot help---there is no clean channel to
switch to---matching the single-modality benchmark. Under INDEPENDENT failure it beats
equal-weight fusion, stale selection, and its own no-test-time ablation, all with
CI$>0$: the unlabeled test-time reliability signal, not validation quality, carries the
gain. The negative result is thus not a refutation of reliability gating in general but
a statement that real single-modality detectors violate its enabling condition.

\section{Operating-Boundary Theory}\label{sec:theory}
The simulation instantiates two boundary theorems (proofs in Appendix~\ref{sec:proofs}).

\begin{theorem}[Clean ceiling]\label{thm:clean}
When validation and test are drawn from the same distribution and detectors are
calibrated on validation, no reliability reweighting of test scores improves expected
AUROC over the validation-optimal combination in expectation. Validation quality is a
sufficient statistic; reliability features are redundant.
\end{theorem}

\begin{theorem}[Independent-failure gain]\label{thm:stress}
If at deployment a strict subset of channels loses discriminative signal while the
complement remains informative, and the failure is detectable from unlabeled test-score
disagreement, then a consensus-agreement reweighting achieves strictly lower expected
regret than both equal-weight fusion and stale validation selection.
\end{theorem}

Theorem~\ref{thm:clean} explains the i.i.d.\ null; Theorem~\ref{thm:stress} explains the
INDEPENDENT-regime gain and predicts where reliability routing can help. Single-modality
data fails the antecedent of Theorem~\ref{thm:stress} (channels share inputs, so failure
is not a strict subset)---hence the negative result is theory-consistent.

\section{Pre-registered Multimodal Test}\label{sec:multimodal}
Theorem~\ref{thm:stress} and the simulation make a falsifiable prediction for genuine
multimodal anomaly detection (RGB $+$ 3D), where a modality can fail independently. We
pre-register the test (frozen decision rule, no post-hoc retuning): under independent
per-modality deployment degradation, reliability-gated fusion must beat (H1) equal-weight
fusion, (H2) stale validation selection, and (H3) its no-test-time ablation, each with
paired-bootstrap CI$>0$ over (category $\times$ degradation-seed) units. A CPU positive
control with the correct design passes decisively (H1 $+0.094$, H2 $+0.264$, H3 $+0.097$,
all CI$>0$); the real test runs on MVTec-3D / Real-IAD. A pass yields a scoped
novel-method claim; a failure is reported as the final boundary, with no novel claim.

\section{Honest Scope and Limitations}
The positive baseline and the negative reliability result are established only on
single-modality, feature-vector anomaly tasks; we make no industrial-vision or
time-series claim. The reliability gate we test is one (consensus-agreement)
instantiation; we cannot exclude that a different reliability parameterization helps,
though three regimes and two designs make it unlikely on this data. The mechanism
section is a controlled simulation, not real-data evidence; it motivates rather than
proves the multimodal claim, which is why we pre-register it. The fraud family has
$n{=}1$ and is excluded from family-level conclusions.

\section{Conclusion}
On a broad anomaly-detection benchmark, a simple rank-normalized stack is a strong,
under-reported universal baseline, and reliability-aware routing---despite its intuitive
appeal---does not improve on it. We trace this to a single structural condition:
reliability gating requires an independently failing channel, which single-modality
detection cannot provide. The honest contribution is a map of where the idea helps and
where it cannot, plus a pre-registered test for the one regime that satisfies the
condition.

\appendix
\section{Proofs}\label{sec:proofs}
\begin{proof}[Proof sketch of Theorem~\ref{thm:clean}]
If $(S_{\text{val}},y_{\text{val}})$ and $(S_{\text{test}},y_{\text{test}})$ are i.i.d.,
the conditional score distributions $p(s\mid y)$ are identical across splits. The AUROC
of any fixed scalarization $w^\top r(x)$ is a functional of these conditionals; hence the
weights maximizing validation AUROC maximize expected test AUROC. A reliability
reweighting that depends only on the (label-free) marginal $p(s)$ cannot improve on the
validation-optimal $w$ in expectation, since the marginal carries no information about
$p(s\mid y)$ beyond what validation already used. $\square$
\end{proof}
\begin{proof}[Proof sketch of Theorem~\ref{thm:stress}]
Partition channels into informative $\mathcal{I}$ and failed $\mathcal{F}$ at test.
Failed channels have test scores independent of $y$, so they (i) add zero-mean noise to
equal-weight fusion, raising regret, and (ii) are detectable by low rank-correlation with
the $\mathcal{I}$-consensus. A weighting $w_m \propto q_m\,a_m$ with $a_m$ the
consensus-agreement drives $w_{\mathcal{F}}\to 0$, recovering the fusion restricted to
$\mathcal{I}$, whose expected AUROC strictly exceeds both the diluted equal-weight fusion
and any single (possibly stale-selected) channel when $|\mathcal{I}|\ge 2$. $\square$
\end{proof}

\end{document}
